\section{Introduction}
\label{sec:intro}

Large language models are increasingly deployed in production systems that process continuously evolving input streams---from customer support queries to news summarization pipelines. When the input distribution shifts, model predictions can silently degrade without any explicit error signal~\cite{khaki2023uncovering}. Detecting such distribution drift early is critical for maintaining system reliability.

Standard drift detectors monitor low-order statistics of embedding distributions: centroid shift tracks the mean, covariance shift tracks the second moment, and Maximum Mean Discrepancy (MMD) measures distributional distance in a reproducing kernel Hilbert space~\cite{khaki2023uncovering, sodar2023detecting}. These methods are effective when drift manifests as shifts in location or spread, but they may miss subtler forms of drift that reorganize the \emph{geometric structure} of the embedding manifold without changing its center of mass. For example, if user queries shift from a diverse cluster to a ring-like distribution around a fixed topic---avoiding certain sub-topics while converging on others---the centroid remains stable while the topology changes fundamentally.

Topological Data Analysis (TDA), and persistent homology in particular, captures multi-scale geometric features of point clouds that are invariant to coordinate transformations~\cite{tralie2019ripser, tauzin2020giotto}. Recent work has applied TDA to characterize LLM representations~\cite{gardinazzi2024persistent}, detect LLM-generated text~\cite{wei2025shortphd}, and monitor concept drift in image streams~\cite{basterrech2024unsupervised}. \citet{yang2026matheval} proposed a general framework for topological drift detection in LLM evaluation, using persistent homology with formal stability guarantees to track structural changes in model capability landscapes, and identified embedding-stream monitoring as a natural high-dimensional instantiation. However, no prior work systematically compares TDA-based drift detectors against classical baselines on real LLM embedding streams across multiple drift types, datasets, and embedding models.

We address this gap with the first comprehensive comparative study of TDA versus classical drift detection on sentence embedding streams. Our experimental setup significantly extends prior work:

\begin{itemize}[leftmargin=*,itemsep=0pt,topsep=0pt]
    \item \textbf{Multiple datasets and encoders}: We evaluate on \agnews (4 topics, well-separated) and \twentyng (6 categories with closely-related pairs) using two embedding models: \minilm (384-d) and \bertbase (768-d).
    \item \textbf{Harder drift scenarios}: Beyond standard topic drift, we introduce \emph{subtopic reweighting} (shifting proportions within a topic while preserving the overall centroid) and \emph{style perturbation} (rotating low-dimensional subspace structure while preserving centroid and covariance). These centroid-preserving drifts on real data bridge the gap between trivial topic drifts and synthetic topology-only experiments.
    \item \textbf{Expanded baselines}: We add energy distance, a classifier-based two-sample test (logistic regression), persistence landscapes, and sliced Wasserstein distance on persistence diagrams---20+ methods total spanning geometry-aware statistics and stronger TDA features.
    \item \textbf{Comprehensive ablations}: We systematically study the effect of window size (50--400 samples), TDA subsample size (40--160 points), and PCA dimensionality reduction before persistent homology (20, 50, 100 dimensions vs.\ raw), with confidence intervals across 5 random seeds.
    \item \textbf{Methodological fixes}: We use a proper calibration/evaluation split for FPR estimation, report end-to-end per-window cost (including persistent homology computation), and evaluate synthetic experiments with AUC rather than scale-dependent separability ratios.
\end{itemize}

The rest of this paper is organized as follows. \Secref{sec:related} reviews related work on drift detection and TDA for NLP. \Secref{sec:method} describes our methodology. \Secref{sec:results} presents results on real and synthetic data plus ablation studies. \Secref{sec:discussion} discusses implications and limitations. \Secref{sec:conclusion} concludes.
